
\Problem{Ganadores}

\Statement

Suponiendo la declaración de la tupla \verb|Player| siguiente
\begin{verbatim}
struct Player { 
    string name; 
    int score;
};
\end{verbatim}
haz una función con la cabecera
\begin{verbatim}
vector<string> winners(vector<Player>::iterator begin,
                       vector<Player>::iterator end);
\end{verbatim}
que devuelva un vector con los jugadores con la puntuación máxima.
Ten en cuenta que puede haber varios.

\Input

La función recibe dos iteradores indicando principio y final de una subsecuencia de un vector de jugadores. Esta subsecuencia puede tener cualquier
longitud, incluso estar vacía. La subsecuencia también podría ser el vector entero.

\Output

La función devuelve un vector con los nombres de los jugadores que tienen
a la vez la puntuación máxima en el mismo orden que se encuentran en el vector
original.

\Observation

En los archivos públicos (icono del gatito) encontrarás: \verb|main.cc|, \verb|player.hh| (declaración de la tupla \verb|Player|), \verb|guanyadors.cc| (donde debes implementar la función). También hay un \verb|Makefile| y el directorio \verb|.vscode| que tiene la configuración para compilar y depurar con VSCode.

Debes implementar \verb|winners| en el archivo \verb|guanyadors.cc|, compilar (está preparado para poder compilar y depurar con VSCode), y enviar solo el archivo \verb|guanyadors.cc|.
\SampleTwoCol

